Newspapers shape public discourse, but their volume makes it impossible for a person to read a whole year of coverage and describe what was reported and in what tone. This project builds and evaluates a complete, automated pipeline that turns page images of a printed newspaper into topic and sentiment measurements, and uses it to study a full year (2025) of The Indian Express, Jaipur edition.

The pipeline has five stages. Page images of the front page (every day) and the editorial page (Monday to Saturday) are collected from the e-paper, giving 663 pages. Each page is preprocessed (deskewing, photograph removal, Sauvola binarisation, removal of printed rules) and segmented into articles by a rule-based layout analyser built on Tesseract line detection, which keeps headlines, decks and body columns in reading order and discards mastheads, advertisements and captions. Every article is then OCR'd, cleaned and filtered to at least 50 words. The result is 5,648 articles and about 1.44 million words, with a mean OCR confidence of 89.8.

Sentiment is measured with three methods of different kinds: the VADER lexicon, a RoBERTa transformer applied sentence by sentence, and the large language model Claude Haiku, which judges the whole article against written definitions and returns a label, a score and a one-line reason. Topics are found with BERTopic (MiniLM embeddings, UMAP, HDBSCAN and class-based TF-IDF), which was chosen in a preliminary comparison with LSA and LDA on a corpus of 329 OCR'd newspaper pages (Cv of 0.69 against 0.42 and 0.39). On the full year it finds 18 topics.

The three sentiment methods describe very different newspapers. Claude Haiku labels 45\% of the articles negative, VADER labels 48\% positive and the transformer labels 92\% neutral; all three agree on only 14\% of articles. Against 48 hand-labelled articles, Claude Haiku is clearly closest to a human reader (Cohen's $\kappa = 0.45$, macro-F1 $= 0.63$), the transformer is no better than always answering ``neutral'' ($\kappa = 0.07$), and VADER has the lowest accuracy (40\%). An analysis of VADER's errors shows that legal, political and diplomatic vocabulary such as ``supreme'', ``justice'', ``party'' and ``security'' is scored as praise.

Although absolute sentiment levels depend on the method, the relative patterns are stable. Space, cricket and cultural coverage is the most positive; crime, courts, Gaza and Iran, Bangladesh and India--Pakistan coverage is the most negative. Monthly sentiment shows no trend over the year but event-driven dips that coincide with surges in the matching topics: the Pahalgam attack and the India--Pakistan hostilities in April and May, the Israel--Iran war in June, the US tariffs in August and the Red Fort car blast in November. Under Claude Haiku the front page is more negative than the editorial page, while the transformer shows the opposite direction. The work provides a reproducible, resumable tool chain and a documented comparison of three sentiment paradigms on long-form journalistic text.

\vspace{2mm}
\noindent\textbf{Keywords:} Topic Modeling, Sentiment Analysis, BERTopic, VADER, Large Language Models, Newspaper NLP, OCR, Layout Analysis, Text Mining
